% Options for packages loaded elsewhere
\PassOptionsToPackage{unicode}{hyperref}
\PassOptionsToPackage{hyphens}{url}
\documentclass[
]{article}
\usepackage{xcolor}
\usepackage[margin=1in]{geometry}
\usepackage{amsmath,amssymb}
\setcounter{secnumdepth}{-\maxdimen} % remove section numbering
\usepackage{iftex}
\ifPDFTeX
  \usepackage[T1]{fontenc}
  \usepackage[utf8]{inputenc}
  \usepackage{textcomp} % provide euro and other symbols
\else % if luatex or xetex
  \usepackage{unicode-math} % this also loads fontspec
  \defaultfontfeatures{Scale=MatchLowercase}
  \defaultfontfeatures[\rmfamily]{Ligatures=TeX,Scale=1}
\fi
\usepackage{lmodern}
\ifPDFTeX\else
  % xetex/luatex font selection
\fi
% Use upquote if available, for straight quotes in verbatim environments
\IfFileExists{upquote.sty}{\usepackage{upquote}}{}
\IfFileExists{microtype.sty}{% use microtype if available
  \usepackage[]{microtype}
  \UseMicrotypeSet[protrusion]{basicmath} % disable protrusion for tt fonts
}{}
\makeatletter
\@ifundefined{KOMAClassName}{% if non-KOMA class
  \IfFileExists{parskip.sty}{%
    \usepackage{parskip}
  }{% else
    \setlength{\parindent}{0pt}
    \setlength{\parskip}{6pt plus 2pt minus 1pt}}
}{% if KOMA class
  \KOMAoptions{parskip=half}}
\makeatother
\usepackage{longtable,booktabs,array}
\usepackage{calc} % for calculating minipage widths
% Correct order of tables after \paragraph or \subparagraph
\usepackage{etoolbox}
\makeatletter
\patchcmd\longtable{\par}{\if@noskipsec\mbox{}\fi\par}{}{}
\makeatother
% Allow footnotes in longtable head/foot
\IfFileExists{footnotehyper.sty}{\usepackage{footnotehyper}}{\usepackage{footnote}}
\makesavenoteenv{longtable}
\usepackage{graphicx}
\makeatletter
\newsavebox\pandoc@box
\newcommand*\pandocbounded[1]{% scales image to fit in text height/width
  \sbox\pandoc@box{#1}%
  \Gscale@div\@tempa{\textheight}{\dimexpr\ht\pandoc@box+\dp\pandoc@box\relax}%
  \Gscale@div\@tempb{\linewidth}{\wd\pandoc@box}%
  \ifdim\@tempb\p@<\@tempa\p@\let\@tempa\@tempb\fi% select the smaller of both
  \ifdim\@tempa\p@<\p@\scalebox{\@tempa}{\usebox\pandoc@box}%
  \else\usebox{\pandoc@box}%
  \fi%
}
% Set default figure placement to htbp
\def\fps@figure{htbp}
\makeatother
\setlength{\emergencystretch}{3em} % prevent overfull lines
\providecommand{\tightlist}{%
  \setlength{\itemsep}{0pt}\setlength{\parskip}{0pt}}
\usepackage{graphicx}
\usepackage{float}
\usepackage{bookmark}
\IfFileExists{xurl.sty}{\usepackage{xurl}}{} % add URL line breaks if available
\urlstyle{same}
\hypersetup{
  hidelinks,
  pdfcreator={LaTeX via pandoc}}

\author{}
\date{\vspace{-2.5em}}

\begin{document}

\section{Question}\label{question}

Las ruinas de Chichén Itzá en Yucatán es uno de los sitios turísticos
más representativos de México. En la tabla se relaciona la cantidad de
personas que ingresó cada día durante una semana, según el tipo de
entrada que pagó.

Se pagan 18 pesos mexicanos de entrada y 7 más si se realiza reserva.

\begin{longtable}[]{@{}
  >{\raggedright\arraybackslash}p{(\linewidth - 14\tabcolsep) * \real{0.2105}}
  >{\centering\arraybackslash}p{(\linewidth - 14\tabcolsep) * \real{0.0921}}
  >{\centering\arraybackslash}p{(\linewidth - 14\tabcolsep) * \real{0.1053}}
  >{\centering\arraybackslash}p{(\linewidth - 14\tabcolsep) * \real{0.1447}}
  >{\centering\arraybackslash}p{(\linewidth - 14\tabcolsep) * \real{0.1053}}
  >{\centering\arraybackslash}p{(\linewidth - 14\tabcolsep) * \real{0.1184}}
  >{\centering\arraybackslash}p{(\linewidth - 14\tabcolsep) * \real{0.1053}}
  >{\centering\arraybackslash}p{(\linewidth - 14\tabcolsep) * \real{0.1184}}@{}}
\caption{Cantidad de personas que ingresaron}\tabularnewline
\toprule\noalign{}
\begin{minipage}[b]{\linewidth}\raggedright
Tipo de entrada
\end{minipage} & \begin{minipage}[b]{\linewidth}\centering
Lunes
\end{minipage} & \begin{minipage}[b]{\linewidth}\centering
Martes
\end{minipage} & \begin{minipage}[b]{\linewidth}\centering
Miércoles
\end{minipage} & \begin{minipage}[b]{\linewidth}\centering
Jueves
\end{minipage} & \begin{minipage}[b]{\linewidth}\centering
Viernes
\end{minipage} & \begin{minipage}[b]{\linewidth}\centering
Sábado
\end{minipage} & \begin{minipage}[b]{\linewidth}\centering
Domingo
\end{minipage} \\
\midrule\noalign{}
\endfirsthead
\toprule\noalign{}
\begin{minipage}[b]{\linewidth}\raggedright
Tipo de entrada
\end{minipage} & \begin{minipage}[b]{\linewidth}\centering
Lunes
\end{minipage} & \begin{minipage}[b]{\linewidth}\centering
Martes
\end{minipage} & \begin{minipage}[b]{\linewidth}\centering
Miércoles
\end{minipage} & \begin{minipage}[b]{\linewidth}\centering
Jueves
\end{minipage} & \begin{minipage}[b]{\linewidth}\centering
Viernes
\end{minipage} & \begin{minipage}[b]{\linewidth}\centering
Sábado
\end{minipage} & \begin{minipage}[b]{\linewidth}\centering
Domingo
\end{minipage} \\
\midrule\noalign{}
\endhead
\bottomrule\noalign{}
\endlastfoot
Sin reserva & 400 & 400 & 500 & 600 & 500 & 300 & 800 \\
Con reserva & 700 & 900 & 300 & 700 & 400 & 600 & 600 \\
\end{longtable}

El administrador diseñó el siguiente proceso para analizar la
recaudación:

\includegraphics[width=1\linewidth,height=\textheight,keepaspectratio]{diagrama_flujo.png}

¿Cuál de las siguientes preguntas \textbf{NO} se puede resolver
únicamente con la información obtenida mediante este proceso?

\subsection{Answerlist}\label{answerlist}

\begin{itemize}
\tightlist
\item
  ¿Cuál fue la ganancia total de la ruinas de Chichén Itzá durante la
  semana?
\item
  ¿Cuál fue el total de dinero recogido por ambos casos durante la
  semana?
\item
  ¿En cuál de los dos casos se recogió más dinero durante la semana?
\item
  ¿En cuál día de la semana se recogió más dinero?
\end{itemize}

\section{Solution}\label{solution}

Para resolver este problema, debemos analizar qué información
proporciona el proceso descrito:

\textbf{Información que SÍ se puede obtener del proceso:}

\begin{enumerate}
\def\labelenumi{\arabic{enumi}.}
\item
  \textbf{Recaudación por tipo de entrada}: El proceso calcula cuánto
  dinero se recogió de personas con reserva (cantidad × 25 pesos
  mexicanos) y sin reserva (cantidad × 18 pesos mexicanos).
\item
  \textbf{Comparación entre casos}: El paso 3 permite comparar los
  resultados, determinando cuál tipo de entrada generó más ingresos.
\item
  \textbf{Total recaudado}: Sumando ambos resultados se obtiene el total
  de dinero recogido.
\end{enumerate}

\textbf{Información que NO se puede obtener:}

\begin{itemize}
\item
  \textbf{Ganancia total}: La ganancia requiere conocer los costos
  operativos de la ruinas de Chichén Itzá (personal, mantenimiento,
  servicios, etc.), información que no está incluida en el proceso.
\item
  \textbf{Día específico}: El proceso trabaja con totales de la semana,
  no con datos diarios.
\end{itemize}

\textbf{Respuesta correcta}: La pregunta sobre la \textbf{ganancia
total} no se puede resolver porque el proceso solo calcula recaudación
(ingresos), no ganancia (ingresos - costos).

\subsection{Answerlist}\label{answerlist-1}

\begin{itemize}
\tightlist
\item
  \textbf{Verdadero}. Esta pregunta sobre la ganancia NO se puede
  resolver porque el proceso solo calcula recaudación, no ganancia (que
  requiere conocer los costos).
\item
  \textbf{Falso}. Esta pregunta SÍ se puede resolver con la información
  del proceso.
\item
  \textbf{Falso}. Esta pregunta SÍ se puede resolver con la información
  del proceso.
\item
  \textbf{Falso}. Esta pregunta SÍ se puede resolver con la información
  del proceso.
\end{itemize}

\section{Meta-information}\label{meta-information}

exname:
proceso\_recaudacion\_sitio\_turistico\_numerico\_variacional\_argumentacion\_n2\_v1
extype: schoice exsolution: 1000 exshuffle: TRUE exsection:
Numerico-Variacional/Argumentacion

\end{document}
